\chapter*{Introduction générale}
\addcontentsline{toc}{chapter}{Introduction générale}

Dans le cadre du projet porté par l'association Tockem, un ensemble d'activités visant à rendre l'accès à l'eau potable durable pour les populations de la localité a été mis en œuvre. Ces activités incluent notamment la réhabilitation et la création de réseaux d'adduction d'eau potable, ainsi que la sensibilisation des populations aux bonnes pratiques de consommation et d'utilisation des ressources mises en place.

Cependant, la gestion d'un réseau d'adduction d'eau potable nécessite souvent la réalisation simultanée de plusieurs tâches, telles que la maintenance des installations, la relève des index de compteurs, l'enregistrement et le branchement de nouveaux abonnés, le maintien de la disponibilité de l'eau, la facturation, le recouvrement et la gestion des finances. Ces activités exigent une vision globale et permanente du système.

Actuellement, la commune de Fokoué ne dispose pas de moyens efficaces pour mener à bien ces tâches, ce qui entraîne des difficultés, notamment pour présenter des chiffres clairs sur l'état des réseaux ou pour facturer dans des délais raisonnables, compte tenu des avancées technologiques.

À cet effet, le logiciel \textbf{Tockem SPE} a été développé pour permettre à la commune de Fokoué de gérer ses installations et ses abonnés de manière simple et efficace. Avec Tockem SPE, la commune pourra effectuer des relevés de compteurs cohérents, réaliser des facturations rapides tenant compte des impayés et/ou des avances, et obtenir une vision globale et détaillée du système en temps réel grâce à des tableaux de bord et des graphiques.

Ce document constitue une documentation utilisateur présentant les fonctionnalités du logiciel Tockem SPE, de son installation aux manipulations nécessaires pour réaliser l'ensemble des opérations possibles.
